\section{Método}

La estrategia metodológica corresponde a un estudio de casos múltiples con unidades incrustadas estructurado en tres fases (Semanas 1 a 13). El diseño explica la toma de decisión B2B y la permanencia laboral mediante Análisis Cualitativo Comparado (QCA), viñetas pareadas y rastreo de procesos.

\subsection{Fase 1: Desk Research y Marco Operativo (Semanas 1--2)}
Análisis sistemático de fuentes secundarias, literatura sociológica y registros históricos de JxR para caracterizar la línea base y seleccionar la muestra.

\begin{itemize}
    \item \textbf{Registros internos JxR:} Reconstrucción de la permanencia observada a 30, 60 y 90 días desagregada por género y sistematización de motivos de interrupción en colocaciones previas.
    \item \textbf{Marco sociológico e interseccional:} Integración de teoría sobre estigma, desistimiento delictual, roles de decisión organizacional, arraigo laboral e interseccionalidad (género, antecedente penal y labores de cuidado).
    \item \textbf{Acceso a datos de Gendarmería de Chile:} Ante la ausencia de microdatos abiertos individuales por la Ley N° 19.628 \parencite{Ley19628SobreProteccionDeLaVidaPrivada}, la información de línea base se construirá a partir de reportes consolidados en Power BI y compendios estadísticos oficiales ya disponibles. En paralelo, se gestionarán solicitudes vía Portal de Transparencia (Ley N° 20.285) \parencite{Ley20285SobreAccesoALaInformacionPublica} para obtener microdatos anonimizados complementarios (.csv, .xlsx) sin condicionar los plazos de la Fase 1.
    \item \textbf{Productos de fase:} Línea base de permanencia, matriz de selección de casos por desenlace, libro de códigos inicial, viñetas pareadas y listado de 10 barreras.
\end{itemize}

\subsection{Fase 2: Levantamiento de Casos Múltiples (Semanas 3--7)}
Levantamiento cualitativo sobre 10 casos de empresas seleccionadas por su desenlace en la cadena de decisión (Tabla~\ref{tab:casos}), completando 29 instancias de entrevista (Tabla~\ref{tab:muestra}).

\begin{table}[h]
\centering
\caption{Selección de Casos de Empresas por Desenlace (N=10)}
\label{tab:casos}
\vspace{0.2cm}
\small
\begin{tabular}{lcl}
\toprule
\textbf{Desenlace del Caso} & \textbf{N° Casos} & \textbf{Función Analítica} \\
\midrule
Contrataron y sostuvieron ($>3$ meses) & 3 & Identificar condiciones de permanencia \\
Contrataron y se interrumpió & 3 & Identificar puntos de quiebre operativos \\
Evaluaron y no avanzaron & 3 & Analizar barreras efectivas de adopción \\
Nunca lo consideraron & 1 & Evaluar etapa de desconocimiento inicial \\
\bottomrule
\end{tabular}
\end{table}

\begin{table}[h]
\centering
\caption{Estructura de la Muestra Cualitativa (29 Instancias)}
\label{tab:muestra}
\vspace{0.2cm}
\small
\begin{tabular}{lcl}
\toprule
\textbf{Unidad de Levantamiento} & \textbf{Instancias} & \textbf{Composición} \\
\midrule
Casos de Empresa & 18 & 2 tríadas (Dirección, RRHH, Jefatura), 4 díadas, 4 únicos \\
Mujeres en Reinserción & 8 & 4 con permanencia $>3$ meses, 4 con interrupción (3 en díada) \\
Ecosistema JxR & 1 & Sesión grupal con 6 integrantes de la red \\
Actores Públicos & 2 & Entrevistas sobre fomento y política pública \\
\bottomrule
\end{tabular}
\end{table}

\textbf{Técnicas e Instrumentos:}
\begin{enumerate}
    \item \textbf{Viñetas Pareadas:} Evaluación intra-sujeto de dos perfiles idénticos alternando solo el sexo para aislar el sesgo de género.
    \item \textbf{Priorización de Barreras:} Ordenamiento cualitativo mediante tarjetas de las 3 principales barreras entre 10 opciones.
    \item \textbf{Evaluación Monádica JxR:} Revisión rotada de fichas de Bolsa Laboral, OTEC y CET para identificar ajustes operacionales.
    \item \textbf{Díadas Empresa--Trabajadora:} Entrevistas espejadas a empresa y trabajadora en 3 casos para contrastar el proceso laboral.
\end{enumerate}

\textbf{Confidencialidad y Ética:} Validación por comité de ética externo, consentimientos informados y protocolo de resguardo de datos sensibles.

\subsection{Fase 3: Análisis Causal, Tipología y Propuesta (Semanas 8--13)}
\begin{itemize}
    \item \textbf{Análisis Cualitativo Comparado (QCA):} Modelación de la adopción B2B mediante tabla de verdad evaluando 4 condiciones (patrocinador en dirección, compatibilidad de cuidados, acompañamiento externo y participación de jefatura). Se declara diversidad limitada analizando solo configuraciones empíricas.
    \item \textbf{Comparación Emparejada y Rastreo de Procesos:} Análisis comparativo de los 6 casos que contrataron (sostuvieron vs. interrumpieron) emparejados por rubro y tamaño para determinar los mecanismos de permanencia.
    \item \textbf{Tipología B2B y Ajustes Programáticos:} Elaboración de perfiles empresariales y recomendaciones concretas para Bolsa Laboral, OTEC y CET.
\end{itemize}